\chapter{引言：虚拟化的形式要求与工程维度}\label{sec:intro}

虚拟化试图解决一个看似矛盾的问题：让一台客户机（guest）获得整台机器的使用接口，却始终不把对真实机器资源的最终控制权交给它\cite{source1}。本篇综述围绕这一中心问题，沿着一条从纯软件到硬件辅助再到嵌套组合的路径展开，考察一台虚拟机器如何先在一个普通进程里被完整构造出来，哪些高频操作可以改由 CPU、内存管理单元、中断控制器和 IOMMU 承担，以及当把 hypervisor 本身也当作客户运行时，资源映射、事件归属和管理策略如何跨层组合、又如何减少层间往返。全文以 CPU、内存、中断、设备四个维度作为统一的分析框架，对每一维度都回答同一组问题：资源与状态由谁控制、关键路径怎样执行、何时需要虚拟化管理层介入、介入的代价何在。第~\ref{sec:software}~章以 QEMU 的全系统模拟说明软件实现，第~\ref{sec:hardware}~章以 Armv8.6-A 与 Linux KVM 说明硬件辅助，第~\ref{sec:nested}~章与第~\ref{sec:hyperturtle}~章分别以 Turtles 与 HyperTurtle 说明嵌套虚拟化，前者解决多层语义如何复用单层硬件支持，后者解决跨层往返如何削减；各节在同一四维结构下逐层对照。

\section{背景与范围}\label{sec:intro-bg}

对虚拟化最直接的需求来自多系统共存：在同一批物理资源上同时运行互不信任、甚至互不兼容的软件栈，需要把资源复用、故障隔离和接口兼容统一在一个机制里。隔离要求一个程序无法触及未分配给它的资源，兼容要求客户操作系统及其应用程序无需修改即可运行，复用则要求多个客户共享底层硬件时仍然各得其所。为此引入的管理层称作虚拟化管理器（hypervisor）或虚拟机监视器（VMM），它在一个真实机器上构造出若干台虚拟机（virtual machine），每台虚拟机对外呈现一套独立的指令集、寄存器、地址空间、中断与设备接口；客户看到的执行单元是虚拟 CPU（vCPU），其物理载体是真实处理器上受 VMM 调度的执行上下文。

本节讨论的范围限定在系统级虚拟机，即能够承载未修改操作系统的整机虚拟化，不涉及进程级的语言虚拟机。需要先澄清一个术语上的容易误解：软件模拟与硬件辅助并不是互斥的产品分类，而是实现同一目标的两条路径，且常常并存于一个系统之中。QEMU 既可以独立提供完整的机器模型，以动态翻译的方式执行客户代码，也可以与 KVM 协作，把 CPU 执行路径交给硬件辅助机制，继续承担内存后备、设备模型与整机编排等职责\cite{source2,source3,source5}。因此本综述把软件与硬件两部分并置，意在比较同一维度上两种机制的取舍，而不是把二者对立起来。

\section{Popek--Goldberg 定义与充分条件}\label{sec:intro-pg}

对“什么是可虚拟化的机器”这一问题的形式化回答，来自 Popek 与 Goldberg 1974 年的论文。他们先用三条性质给出 VMM 的定义，再在简化机器模型上证明一个可判定的充分条件\cite{source1}。这三条性质属于 VMM 的定义本身，而不是三个可以分别取舍的定理：

\begin{enumerate}[label=\arabic*.,leftmargin=2.5em,itemsep=2pt]
\item \textbf{等价性（equivalence）}：任何程序在 VMM 之下的执行效果与它直接在原机器上运行不可区分，但明确允许两类例外：
  \begin{itemize}[label=--,leftmargin=1.5em,itemsep=1pt]
  \item 一方面，允许因可用资源数量不同造成的差异，以容纳 VMM 分配不同大小内存的情形；
  \item 另一方面，允许因时间依赖（timing dependencies）造成的差异，因为中间多了一层软件，且同一硬件上还有其他虚拟机并发存在。
  \end{itemize}
\item \textbf{资源控制（resource control）}：VMM 对内存、外设等资源握有完全控制权。一方面，运行于其下的程序不可能访问任何未显式分配的资源；另一方面，VMM 在规定条件下能够重新取得对已分配资源的控制。
\item \textbf{效率（efficiency）}：统计上占主导地位的那一子集客户指令必须直接由真实处理器执行，VMM 完全不介入；正是这一条把传统的解释器和纯软件模拟器排除在虚拟机的定义之外。
\end{enumerate}

充分条件的关键在于区分敏感指令（sensitive instruction）与特权指令（privileged instruction）这两个不同概念：

\begin{enumerate}[label=\arabic*.,leftmargin=2.5em,itemsep=2pt]
\item \textbf{特权指令（privileged instruction）}：机器的固有属性，与虚拟化过程无关——在用户态执行会陷入、在管态执行则不陷入；是否特权只取决于机器本身，从机器的操作原理即可判定。
\item \textbf{敏感指令（sensitive instruction）}：执行效果受系统状态影响的指令，分两类：
  \begin{itemize}[label=--,leftmargin=1.5em,itemsep=1pt]
  \item \textbf{控制敏感（control sensitive）}：试图改变可用内存资源数量或处理器模式，且这种改变不经过内存 trap 序列。
  \item \textbf{行为敏感（behavior sensitive）}：执行效果依赖于它在真实内存中的位置或当前模式。位置敏感（location sensitive）的例子是取得物理地址，模式敏感（mode sensitive）的例子是访问上一执行模式的指令空间。
  \end{itemize}
\end{enumerate}

敏感与特权并非同义词：一个机器即使具备管态与用户态两种模式，也未必满足“敏感集合包含于特权集合”；若存在某条敏感指令在用户态不陷入，仅仅把操作系统降权运行就无法保证等价性与资源控制同时成立。正是在这两个集合的关系上，论文给出其核心结论：

\begin{theorem}[Popek--Goldberg, 1974]\label{thm:pg1}
对一台满足论文机器模型假设的第三代式机器，若其敏感指令集合是特权指令集合的子集（sensitive $\subseteq$ privileged），则可以为该机器构造出满足等价性、资源控制与效率三条性质的 VMM。
\end{theorem}

需要强调，这是一个充分条件而非必要条件：不满足该条件的机器未必不可虚拟化。trap-and-emulate 是对这一条件的直观解释——敏感指令在用户态自然陷入，由 VMM 代为模拟其语义——但它描述的是该条件下一种可行的实现方式，并不能外推为所有现代实现的必要条件。

这一形式结果有两个边界必须说清：

\begin{enumerate}[label=\arabic*.,leftmargin=2.5em,itemsep=2pt]
\item \textbf{机器模型的范围}。定理依赖的机器模型——论文所称“第三代式机器”，即 IBM 360/370、Honeywell 6000、DEC PDP-10 这类具备管态/用户态双模式与内存重定位的经典计算机的简化模型——把机器状态记为四元组 $S = (E, M, P, R)$：$E$ 为可执行存储（容量 $q$ 的线性内存），$M$ 为处理器模式（$M \in \{s, u\}$，管态或用户态），$P$ 为程序计数器，$R = (l, b)$ 为 relocation-bounds 寄存器，即把访存地址 $a$ 重定位为 $a + l$ 并检查 $a \le b$ 的重定位界限。模型还显式假设 I/O 指令与中断不存在，尽管论文指出它们可作为扩展加入\cite[§2]{source1}。因此该模型覆盖的是处理器状态与一种简化的重定位内存，没有把完整中断与设备行为纳入形式证明。
\item \textbf{递归虚拟化的前提}。论文第 7 节讨论递归虚拟化——VMM 能否在其自身上再运行一份 VMM——并给出定理 2：
\begin{theorem}[Popek--Goldberg, 1974]\label{thm:pg2}
若一台第三代式机器满足以下两个条件，则它是递归可虚拟化的：
\begin{enumerate}[label=(\alph*),leftmargin=2.5em,itemsep=1pt]
\item 该机器本身可虚拟化；
\item 可以为它构造出一个无时间依赖（无 timing dependencies）的 VMM。
\end{enumerate}
\end{theorem}
递归条件不承诺固定不变的低开销：每增加一层 VMM 都要消耗资源，实际实现还可能增加调度、状态切换与陷入处理的成本；第~\ref{sec:nested}~章与第~\ref{sec:hyperturtle}~章将回到这一点，分别讨论其工程实现与跨层代价，而不是宣称定理在第二层失效。
\end{enumerate}

\section{从形式模型到四维工程框架}\label{sec:intro-dims}

Popek--Goldberg 的形式模型为“可虚拟化”提供了可检验的判据，但它的抽象范围与现代整机虚拟化所面对的对象并不重合。它把内存抽象成一条 relocation-bounds 寄存器驱动的简单重定位，把 I/O 指令与中断整个排除在外；而真实系统里的内存虚拟化涉及多级页表与访问权限，中断虚拟化涉及路由、pending 状态与投递时机，设备虚拟化则要处理 MMIO、DMA 与设备状态。要把理论判据落到工程实现上，就需要把分析对象补全。表 \ref{tab:fourdim} 给出本文采用的四个维度，它用于分解整机虚拟化所需的机制，是后续各节的共同分析框架，而非 Popek--Goldberg 原文提出的分类；中断与设备在物理上虽有交叉，但本文把设备维度留给功能与数据通路，把中断维度留给异步控制流，分别讨论。

\begin{table}[htbp]
\centering
\caption{四维分析框架：各维度的对象、控制目标与共同问题}
\label{tab:fourdim}
\begin{tabularx}{\linewidth}{@{}l>{\raggedright\arraybackslash}X>{\raggedright\arraybackslash}X>{\raggedright\arraybackslash}X@{}}
\toprule
维度 & 客户得到什么 & VMM 必须控制什么 & 各章共同回答的问题 \\
\midrule
CPU & 指令集、寄存器、特权状态与 vCPU & 执行上下文、敏感操作、运行与暂停 & 指令由谁执行？何时需要 VMM 介入？ \\
内存 & 自己的物理地址空间与页表 & 真实内存分配、地址映射与访问权限 & 客户认为的“物理地址”最终指向哪里？ \\
中断 & 定时器及设备产生的异步事件 & 中断路由、pending 状态、目标 vCPU 与投递时机 & 事件如何送到正确的客户，并遵守其屏蔽与优先级规则？ \\
设备 & 可发现、可编程的设备接口 & MMIO、DMA、设备状态与后端资源 & 寄存器访问、数据传输和完成通知分别经过谁？ \\
\bottomrule
\end{tabularx}
\end{table}

最后还需划定一个术语边界。本文把 QEMU 的跨指令集全系统模拟纳入广义的“软件虚拟化”，但必须指出它与 Popek--Goldberg 意义上的 efficiency 不重合：efficiency 要求统计上占主导的客户指令由真实处理器直接执行，而 QEMU 将 RISC-V 指令翻译成中间表示再生成 x86-64 代码执行，二者并非同一指令在真实处理器上的直接执行，因此并不满足这条严格定义\cite{source1,source3}。翻译块（TB）缓存降低的是反复翻译与解码的模拟开销，它提高的是模拟效率，不能把它偷换成“直接执行”。这一区分并非贬低软件路径的价值：软件路径不依赖宿主提供目标架构的虚拟化扩展；在模拟器具有对应 CPU、设备模型及可用宿主后端的前提下，它能在异构平台上运行整套 guest 软件。下一节就从这条软件路径出发，看一台虚拟机器如何在 QEMU 进程内部被逐维度地构造出来。
